\section{模型求解}\label{sec:solving}
我们利用三次真实点球的早期位置反演初始运动，并比较精准推射与大力抽射的飞行特点和预测结果。
\begin{table}[H]\centering\small\caption{四模型体系及实际反演量}\label{tab:models}
\begin{tabular}{lll}\toprule
模型 & 受力或运动假设 & 反演量\\\midrule
M1 & 重力 & 初速度（3项）\\
M2 & 重力、空气阻力 & 初速度、阻力系数（4项）\\
M3 & 恒合力、匀加速度 & 初速度、加速度（6项）\\
M4 & 重力、阻力、马格努斯力 & 初速度、旋转方向、$C_D$、$\alpha$（7项）\\\bottomrule
\end{tabular}\end{table}
M3 的恒合力包含重力及其他力的共同作用。M4 中 $C_D$ 和 $\alpha$ 固定于每次飞行，升力系数 $C_L=\alpha Sp$ 随运动变化。正向参考值为 $C_D^*=\ReferenceCD$、$\alpha^*=\ReferenceCL$，取自文献参考工况\cite{kiratidis2018}。

\subsection{三次真实点球的逆向预测}\label{sec:inverse}
\subsubsection{测量背景、坐标与输入时刻}
Sportec Open DFL 数据将比赛事件与 TRACAB Gen5 多相机光学追踪关联，采样频率为 25 Hz\cite{bassek,dataset}。本文采用 Dejan Ljubicic（J03WMX）、Dawid Igor Kownacki（J03WOH）及 Philipp Klement（J03WR9）的点球。所给 CSV 共保留 39 个源帧，三例分别为 15、12、12 帧。计算沿用发布位置，不另加观测噪声模型。

源水平坐标以球场中心为原点，三例均朝向 $X=52.5$ m。为统一到题目的点球点坐标，采用
\begin{equation}
x=X-41.5,\qquad y=Y,\qquad z=Z_{\rm original},\qquad x_g=11\ \mathrm{m}.
\label{eq:data-coordinates}\end{equation}
名义点球点为水平原点，各次首个实测位置保留自身偏移。发布的 Z 原值按米及球心高度解释；质量、半径和空气参数采用第二章的模型参考设定，数据来源未给出逐球实测质量和半径。

每次反演以所选起始源帧为 $t=0$。J03WMX 使用原序列第 3、4、5 个采样点，J03WOH 与 J03WR9 使用第 2、3、4 个点。三例相对拟合时刻均为 0、0.04、0.08 s，所选起点在原时间轴上分别为 0.08、0.04、0.04 s。表~\ref{tab:penalty-inputs} 列出新的实测起点。此前记录、重复帧和门线后记录均保留在原始 CSV；轨迹对照从所选起点开始，后续采样仅用于比较，窗口截至实测门线穿越。帧间交点采用线性插值，与实测帧分开标记。
\PenaltyInputsTable

\subsubsection{有界反演与选解}
M3 使用三个位置直接解出恒加速度：
\begin{equation}\boldsymbol r(t)=\boldsymbol r_0+\boldsymbol v_0t+\tfrac12\boldsymbol a t^2.\end{equation}
M4 以初速度、旋转方向、$C_D$、$\alpha$ 为变量，最小化两个后续位置的平方误差。确定性六组初值的最小残差对应代表解；各参数范围列于表~\ref{tab:inverse-bounds}。
\InverseBoundsTable
马格努斯力包含乘积 $\alpha\boldsymbol\omega$，仅由位置不能分别唯一确定 $\alpha$ 和自旋大小。为选择代表解，令 $\|\boldsymbol\omega\|=30$ rad/s；三个初速度、两个方向角和两个系数共有七项，而两个后续位置提供六个坐标，因此结果仍受边界与选解条件影响。

\subsection{案例分析与验证}\label{sec:case-validation}
三次点球分为两种射门方式：J03WMX 和 J03WR9 为精准推射，J03WOH 为大力抽射。图~\ref{fig:penalty-overview} 汇总三次实测空间轨迹及模型预测。空心圆表示反演输入，黑线和黑点表示实测位置，彩色线表示四种模型，叉号表示实测或插值得到的球门交点；预测在首次触地时结束。
\begin{figure}[H]\centering\includegraphics[width=\textwidth]{penalty_spatial_overview.pdf}
\caption{三次真实点球的空间轨迹与模型预测。每次反演使用选定起点后 0、0.04、0.08 s 的三个位置。}\label{fig:penalty-overview}\end{figure}

\subsubsection{案例一：精准推射}
J03WMX 和 J03WR9 的实测轨迹均贴近草地飞行。M4 在两个反演起点给出的球速分别为 29.20 和 29.33 m/s。J03WMX 中 M4 到达球门平面，交点误差为 0.360 m；M1、M2 提前触地，M3 在计算时限内未到达球门。J03WR9 中 M1、M2 的交点误差分别为 0.550 和 0.541 m，而 M3、M4 预测为提前触地。

低空轨迹对竖直受力较敏感：重力持续使球下落，竖直方向的马格努斯力则可抬高或压低轨迹。M1 忽略气动力，M2 加入减速阻力，却没有横向升力；M3 用三个点给出恒加速度，能够通过早期位置，但后续加速度随速度变化时，其外推轨迹可能偏离。M4 允许气动力随速度和旋转方向变化；J03WMX 的结果更接近球门交点，而 J03WR9 的预测仍提前触地。具体误差还受到早期位置和拟合参数共同影响。

\subsubsection{案例二：大力抽射}
J03WOH 的 M4 反演起点球速为 32.05 m/s，高于上述两次推射。图~\ref{fig:penalty-J03WOH} 展示空间轨迹及 $x$、$y$、$z$ 坐标随时间的变化；前三个位置用于拟合，后续黑点用于检验预测。
\begin{figure}[H]\centering\includegraphics[width=\textwidth]{penalty_J03WOH.pdf}
\caption{J03WOH 大力抽射的空间轨迹及三个坐标随时间的变化。空心圆为拟合输入，黑点为实测采样。}\label{fig:penalty-J03WOH}\end{figure}
$x(t)$ 体现前进方向的减速，$y(t)$ 与 $z(t)$ 分别体现横向偏转和竖直弯曲。M1、M2、M3、M4 的球门交点误差依次为 0.684、0.657、0.374 和 0.113 m。M1 只考虑重力，M2 进一步计入阻力，两者的横向方向主要由初速度决定；M3 延续三点确定的恒加速度；M4 则通过阻力与马格努斯力描述随速度变化的弯曲，在该次射门中与后续实测更接近。交点误差按模型各自穿越球门平面时的 $(y,z)$ 与实测交点计算。

\paragraph{速度与气动力的物理图像。} 图~\ref{fig:penalty-dynamics} 对比三次射门的 M4 球速和气动力。阻力与相对速度平方成正比；J03WOH 的反演起点速度较高，其初始阻力约 7.58 N，两次推射约为 7.05 和 7.11 N，其中系数取值也会影响大小。随足球减速，阻力逐渐减小。马格努斯力与旋转、相对速度的叉乘有关：旋转方向决定它使球向侧方还是竖直方向弯曲，其大小随速度和旋转条件改变。图中力与速度来自 M4 的反演及数值计算，旋转模长采用 30 rad/s 的参考设定。
\begin{figure}[H]\centering\includegraphics[width=\textwidth]{penalty_m4_dynamics.pdf}
\caption{三次真实点球的 M4 反演球速、阻力和马格努斯力；曲线分别截止于模型首次触地或到达球门。}\label{fig:penalty-dynamics}\end{figure}
